\section{Comunicación}

\subsection{Medida de bytes}
El tamaño de un modelo se calcula como la suma, para cada tensor de su \texttt{state\_dict}, del número de
elementos por los bytes que ocupa cada uno. A partir de esto, en cada ronda se registran:

\begin{itemize}
    \item \textbf{Bajada}: número de clientes seleccionados por el tamaño del modelo global. La bajada no se comprime.
    \item \textbf{Subida}: suma del tamaño de lo que envía cada cliente una vez codificado. De este modo, al activar
    un codec solo se reduce la subida.
\end{itemize}

No se tiene en cuenta el coste de serialización ni el del protocolo de red. El envío se realiza a través de un
codec con dos operaciones: \texttt{encode}, que se ejecuta en el cliente, y \texttt{decode}, que reconstruye los
pesos en el servidor. Ambas reciben como referencia el modelo global de la ronda, que es conocido por los dos
extremos. El codec por defecto no aplica ninguna transformación.

\subsection{Cuantización int8}
En lugar de cuantizar directamente los pesos, cada cliente cuantiza la diferencia $\Delta = w^k - w_t$ respecto al
modelo global, cuyos valores tienen un rango mucho menor y, por tanto, se representan con menos error. Para cada
tensor se calcula una escala $s = \max|\Delta| / 127$ y se envía:
\[
q = \text{clamp}\left(\text{round}\left(\Delta / s\right), -127, 127\right)
\]
como un tensor \texttt{int8}, junto con la escala $s$ en \texttt{float32}. Si todas las diferencias son nulas se
toma $s=1$ para evitar la división entre cero. El servidor reconstruye los pesos como $w_t + q \cdot s$.
\\

Cada parámetro pasa de ocupar 4 bytes a ocupar 1, más 4 bytes por tensor para la escala, por lo que la subida
se reduce aproximadamente a la cuarta parte. Este codec se activa con \texttt{- -set server.codec=int8}.

\subsection{Top-k (extensión)}
Como extensión opcional se implementó una compresión que solo envía las componentes más grandes de la
actualización. Para cada tensor se calcula de nuevo la diferencia $\Delta$ respecto al modelo global y se
seleccionan las $k = \lceil f \cdot n \rceil$ componentes de mayor valor absoluto, siendo $f$ la fracción indicada
en \texttt{server.topk\_fraction} y $n$ el número de elementos del tensor. Se envían sus posiciones en
\texttt{int32} y sus valores en \texttt{float32}, es decir, 8 bytes por componente. En el servidor, las componentes
no recibidas se consideran nulas, de modo que esos pesos se mantienen iguales a los del modelo global.
\\

Al ocupar 8 bytes por componente enviada en lugar de 4, este codec solo reduce la subida si $f < 0.5$. Con el valor
por defecto $f = 0.1$ la subida se reduce a la quinta parte. Se activa con \texttt{--set server.codec=topk}.